% ══════════════════════════════════════════════════════════════════
%  FICHE D'EXERCICES — Chapitre 4 : Les hydrocarbures aromatiques (benzène)
% ══════════════════════════════════════════════════════════════════
\section{Exercices type Baccalauréat}
\textit{Données : $M(\ce{H})=1$, $M(\ce{C})=12$, $M(\ce{O})=16$, $M(\ce{N})=14$,
$M(\ce{Cl})=35{,}5$ \si{\gram\per\mol}.}

\rubrique{Structure \& nomenclature}

\exo{}\diff{1} Représenter la molécule de benzène (formule de Kekulé et
représentation avec le cercle). Indiquer sa formule brute et expliquer le terme
« délocalisation des électrons ».

\exo{}\diff{1} Nommer les composés aromatiques : a)~$\ce{C6H5-CH3}$ ;
b)~$\ce{C6H5-OH}$ ; c)~$\ce{C6H5-NH2}$ ; d)~le 1,4-diméthylbenzène.

\exo{}\diff{2} Le diméthylbenzène (xylène) existe sous trois isomères. Représenter
les trois isomères (ortho, méta, para) et donner leur nom en nomenclature
systématique (positions des substituants).

\rubrique{Réactivité : substitution électrophile}

\exo{}\diff{1} Écrire l'équation de la mononitration du benzène (action de
$\ce{HNO3}$ en présence de $\ce{H2SO4}$). Nommer le produit et préciser le rôle de
l'acide sulfurique.

\exo{}\diff{1} Écrire l'équation de la monochloration du benzène (catalyseur
$\ce{FeCl3}$). Pourquoi parle-t-on de \emph{substitution} et non d'\emph{addition} ?

\exo{}\diff{2} Comparer le comportement du benzène et celui d'un alcène (par exemple
l'éthène) vis-à-vis : a)~de l'eau de brome ; b)~du dibrome en présence de catalyseur.
Conclure sur la stabilité particulière du noyau aromatique.

\exo{}\diff{2} Le benzène réagit avec le chlorométhane $\ce{CH3Cl}$ en présence de
$\ce{AlCl3}$ (réaction de Friedel-Crafts).
\begin{enumerate}[label=\arabic*)]
\item Écrire l'équation de la réaction et nommer le produit.
\item Ce produit peut subir une seconde alkylation : écrire l'équation de formation
d'un diméthylbenzène.
\end{enumerate}

\rubrique{Détermination de formule}

\exo{}\diff{2} Un hydrocarbure aromatique $A$ contient en masse \SI{90.6}{\percent}
de carbone. Sa masse molaire est $M = \SI{106}{\gram\per\mol}$.
\begin{enumerate}[label=\arabic*)]
\item Déterminer la formule brute de $A$.
\item Sachant que $A$ dérive du benzène par substitution, proposer les formules
semi-développées possibles et les nommer.
\end{enumerate}

\exo{}\diff{2} La combustion complète de \SI{3.9}{\gram} de benzène produit du
dioxyde de carbone et de l'eau.
\begin{enumerate}[label=\arabic*)]
\item Écrire l'équation de combustion du benzène.
\item Calculer les masses de $\ce{CO2}$ et d'eau formées.
\end{enumerate}

\rubrique{Problèmes de synthèse — type Baccalauréat}

\sujetbac[d'après Baccalauréat C, Cameroun]
On étudie le toluène (méthylbenzène), de formule $\ce{C6H5-CH3}$.
\begin{enumerate}[label=\arabic*)]
\item Calculer sa masse molaire et son pourcentage massique en carbone.
\item Le toluène subit une mononitration. Écrire l'équation et nommer le(s) produit(s)
possibles.
\item L'oxydation ménagée du groupe méthyle du toluène conduit à l'acide benzoïque
$\ce{C6H5-COOH}$. Écrire la formule de l'acide benzoïque et calculer sa masse molaire.
\item On veut préparer \SI{12.2}{\gram} d'acide benzoïque. Quelle masse de toluène
faut-il, si le rendement de l'oxydation est de \SI{80}{\percent} ?
\end{enumerate}

\sujetbac[Le trinitrotoluène]
La nitration poussée du toluène conduit au 2,4,6-trinitrotoluène (TNT), de formule
$\ce{C6H2(CH3)(NO2)3}$.
\begin{enumerate}[label=\arabic*)]
\item Représenter la molécule de TNT (positions des groupes nitro et méthyle sur le cycle).
\item Calculer la masse molaire du TNT.
\item Écrire l'équation globale de la trinitration du toluène par l'acide nitrique
(on formera également de l'eau).
\end{enumerate}
